% Generated draft: requires mathematical and source review.
\documentclass[11pt]{article}
\usepackage[margin=25mm]{geometry}
\usepackage{fontspec}
\setmainfont{Latin Modern Roman}
\usepackage{amsmath,amssymb,mathtools,unicode-math}
\setmathfont{Latin Modern Math}
\usepackage{enumitem,needspace,xurl,hyperref}
\setlength{\parindent}{0pt}
\setlength{\parskip}{5pt}
\setlength{\emergencystretch}{3em}
\newcommand{\N}{\mathbb{N}}
\newcommand{\Z}{\mathbb{Z}}
\newcommand{\Q}{\mathbb{Q}}
\newcommand{\R}{\mathbb{R}}
\newcommand{\C}{\mathbb{C}}
\newcommand{\F}{\mathbb{F}}
\newcommand{\A}{\mathbb{A}}
\newcommand{\PP}{\mathbb{P}}
\newcommand{\E}{\mathbb{E}}
\newcommand{\eps}{\varepsilon}
\newcommand{\dd}{\,\mathrm{d}}
\begin{document}
\section*{UnsolvedMath ID: 800351}
\textbf{Source problem code:} catalogue entry
\subsection{Definitions and mathematical statement}
Let \(D\) denote the set of admissible parameter pairs for benzels in the cited source. For \((a,b)\in D\), write \(B_{a,b}\) for the \((a,b)\)-benzel, \(I(B_{a,b})\) for its Conway--Lagarias invariant with the source's sign convention, and \(\mathcal T_{103}(B_{a,b})\) for the set of its tilings of type \(103\), using the source's tile and placement conventions. These symbols are placeholders for source definitions: the supplied record does not specify \(D\), the geometry of a benzel, the definition or normalization of the invariant, or the precise meaning of type \(103\).

The question attributed to Propp (2022), Problem 7 (\#3), is whether
\[
\forall (a,b)\in D,\qquad I(B_{a,b})\geq 0
\ \Longrightarrow\ 
\mathcal T_{103}(B_{a,b})\neq\varnothing.
\]
The equality case \(I(B_{a,b})=0\) is explicitly included.

The supplied background states that a strictly negative Conway--Lagarias invariant prevents a benzel from being tiled by bones and right-pointing stones, forcing the corresponding table entries to vanish. It reports observed positive entries when the invariant is nonnegative, while explicitly withholding a general justification. The table and the observations preceding Problem 4 are not supplied. The context associates these tile restrictions with the question about type \(103\), but the exact type encoding and allowed tile orientations are absent. Consequently, the displayed implication records the intended question subject to unresolved source definitions; it is not a fully specified geometric formulation.
\subsection{Short English statement}
Propp asks whether every admissible benzel with nonnegative Conway--Lagarias invariant has a type \(103\) tiling, including when the invariant is zero. The supplied record reports supporting table observations, but omits the definitions needed to specify the regions, invariant, and permitted tilings precisely.
\subsection{Sources}
\begin{itemize}
\item Local UnsolvedMath registry: permanent ID \(800351\), exact source problem code \texttt{catalogue entry}, title \emph{Trimer Covers in the Triangular Grid --- Problem 7}. The supplied record attributes the problem to Propp (2022) and labels it open; this is supplied registry status, not an independently verified literature assessment.
\item Propp, \emph{Trimer Covers in the Triangular Grid: Twenty Mostly Open Problems} (2022), source item Problem 7 (\#3). Supplied source URL: \url{https://arxiv.org/abs/2206.06472}. The record reports access on 2026-10-16 and extraction by \texttt{arxiv\_latex}; these are supplied provenance details, not a claim of independent source access or verification.
\item Supplied upstream revision: \texttt{d371ee917ac9647fbb06cd6fc529deec6f3dda82}. The excerpt refers to observations preceding Problem 4 and to a table, neither of which is included in the supplied record.
\end{itemize}
\end{document}
